\documentclass[12pt,a4paper]{article}
\usepackage[utf8]{inputenc}
\usepackage[french]{babel}
\usepackage[T1]{fontenc}
\usepackage{geometry}
\usepackage{graphicx}
\usepackage{fancyhdr}
\usepackage{titlesec}
\usepackage{tocloft}
\usepackage{listings}
\usepackage{xcolor}
\usepackage{hyperref}
\usepackage{amsmath}
\usepackage{amsfonts}
\usepackage{booktabs}
\usepackage{float}

\geometry{margin=2.5cm}
\pagestyle{fancy}
\fancyhf{}
\fancyhead[L]{LanceJob - Rapport Technique}
\fancyhead[R]{\thepage}

\definecolor{codegreen}{rgb}{0,0.6,0}
\definecolor{codegray}{rgb}{0.5,0.5,0.5}
\definecolor{codepurple}{rgb}{0.58,0,0.82}
\definecolor{backcolour}{rgb}{0.95,0.95,0.92}

\lstdefinestyle{mystyle}{
    backgroundcolor=\color{backcolour},   
    commentstyle=\color{codegreen},
    keywordstyle=\color{magenta},
    numberstyle=\tiny\color{codegray},
    stringstyle=\color{codepurple},
    basicstyle=\ttfamily\footnotesize,
    breakatwhitespace=false,         
    breaklines=true,                 
    captionpos=b,                    
    keepspaces=true,                 
    numbers=left,                    
    numbersep=5pt,                  
    showspaces=false,                
    showstringspaces=false,
    showtabs=false,                  
    tabsize=2
}

\lstset{style=mystyle}

\title{\textbf{LanceJob} \\ Plateforme de Mise en Relation Freelancers-Clients \\ \large Rapport Technique Complet}
\author{Majidi}
\date{\today}

\begin{document}

\maketitle
\newpage

\tableofcontents
\newpage

\section{Introduction}

\subsection{Présentation du Projet}
LanceJob est une plateforme web moderne conçue pour faciliter la mise en relation entre freelancers et clients. Cette application permet aux professionnels indépendants de proposer leurs services et aux clients de trouver les compétences adaptées à leurs besoins.

\subsection{Objectifs}
\begin{itemize}
    \item Créer une interface intuitive pour la gestion des projets freelance
    \item Implémenter un système de matching intelligent entre offres et demandes
    \item Assurer la sécurité des transactions et des données utilisateurs
    \item Optimiser l'expérience utilisateur sur tous les appareils
\end{itemize}

\section{Architecture du Système}

\subsection{Architecture Générale}
Le projet LanceJob adopte une architecture moderne basée sur une séparation claire entre le frontend et le backend :

\begin{itemize}
    \item \textbf{Frontend} : Application React.js avec TypeScript
    \item \textbf{Backend} : API REST développée avec Node.js et Express
    \item \textbf{Base de données} : MongoDB pour la persistance des données
    \item \textbf{Authentification} : JWT (JSON Web Tokens)
\end{itemize}

\subsection{Structure des Dossiers}

\begin{lstlisting}[language=bash, caption=Structure du projet]
lancejob/
├── frontend/
│   ├── public/
│   ├── src/
│   │   ├── components/
│   │   ├── pages/
│   │   ├── services/
│   │   ├── types/
│   │   └── utils/
│   ├── package.json
│   └── tsconfig.json
├── backend/
│   ├── src/
│   │   ├── controllers/
│   │   ├── models/
│   │   ├── routes/
│   │   ├── middleware/
│   │   └── utils/
│   ├── package.json
│   └── server.js
└── reports/
    └── main.tex
\end{lstlisting}

\section{Conception et Modélisation}

\subsection{Modèle de Données}

\subsubsection{Entité Utilisateur}
\begin{itemize}
    \item \textbf{Champs principaux} : ID, nom, email, mot de passe, type (freelancer/client)
    \item \textbf{Profil freelancer} : compétences, portfolio, tarifs, disponibilité
    \item \textbf{Profil client} : entreprise, secteur d'activité, projets publiés
\end{itemize}

\subsubsection{Entité Projet}
\begin{itemize}
    \item \textbf{Informations générales} : titre, description, budget, délais
    \item \textbf{Compétences requises} : technologies, niveau d'expertise
    \item \textbf{Statut} : ouvert, en cours, terminé, annulé
\end{itemize}

\subsubsection{Entité Candidature}
\begin{itemize}
    \item \textbf{Relation} : freelancer ↔ projet
    \item \textbf{Données} : proposition tarifaire, délai proposé, message
    \item \textbf{Statut} : en attente, acceptée, refusée
\end{itemize}

\subsection{Diagrammes UML}

\subsubsection{Diagramme de Classes}
Le système est organisé autour de trois classes principales :
\begin{itemize}
    \item \texttt{User} : gestion des utilisateurs et authentification
    \item \texttt{Project} : gestion des projets et missions
    \item \texttt{Application} : gestion des candidatures
\end{itemize}

\subsubsection{Diagramme de Séquence}
Les interactions principales incluent :
\begin{enumerate}
    \item Inscription/Connexion utilisateur
    \item Publication d'un projet par un client
    \item Candidature d'un freelancer
    \item Sélection et attribution du projet
\end{enumerate}

\section{Technologies Utilisées}

\subsection{Frontend}
\begin{itemize}
    \item \textbf{React.js 18} : Framework JavaScript pour l'interface utilisateur
    \item \textbf{TypeScript} : Typage statique pour une meilleure robustesse
    \item \textbf{Material-UI} : Bibliothèque de composants pour un design moderne
    \item \textbf{React Router} : Gestion de la navigation
    \item \textbf{Axios} : Client HTTP pour les appels API
\end{itemize}

\subsection{Backend}
\begin{itemize}
    \item \textbf{Node.js} : Environnement d'exécution JavaScript côté serveur
    \item \textbf{Express.js} : Framework web minimaliste et flexible
    \item \textbf{MongoDB} : Base de données NoSQL orientée documents
    \item \textbf{Mongoose} : ODM (Object Document Mapping) pour MongoDB
    \item \textbf{bcrypt} : Hachage sécurisé des mots de passe
    \item \textbf{jsonwebtoken} : Implémentation des tokens JWT
\end{itemize}

\subsection{Outils de Développement}
\begin{itemize}
    \item \textbf{Git} : Système de contrôle de version
    \item \textbf{npm} : Gestionnaire de packages
    \item \textbf{ESLint} : Analyse statique du code
    \item \textbf{Prettier} : Formatage automatique du code
\end{itemize}

\section{Implémentation}

\subsection{Configuration du Backend}

\subsubsection{Serveur Express}
\begin{lstlisting}[language=javascript, caption=Configuration du serveur]
const express = require('express');
const cors = require('cors');
const mongoose = require('mongoose');

const app = express();

// Middleware
app.use(cors());
app.use(express.json());

// Routes
app.use('/api/auth', authRoutes);
app.use('/api/projects', projectRoutes);
app.use('/api/users', userRoutes);

// Connexion à MongoDB
mongoose.connect(process.env.MONGODB_URI)
  .then(() => console.log('Connecté à MongoDB'))
  .catch(err => console.error('Erreur de connexion:', err));
\end{lstlisting}

\subsubsection{Modèles Mongoose}
\begin{lstlisting}[language=javascript, caption=Modèle Utilisateur]
const userSchema = new mongoose.Schema({
  name: { type: String, required: true },
  email: { type: String, required: true, unique: true },
  password: { type: String, required: true },
  role: { type: String, enum: ['freelancer', 'client'], required: true },
  profile: {
    skills: [String],
    hourlyRate: Number,
    availability: String,
    portfolio: [String]
  }
}, { timestamps: true });
\end{lstlisting}

\subsection{Configuration du Frontend}

\subsubsection{Composants React}
\begin{lstlisting}[language=typescript, caption=Composant de connexion]
import React, { useState } from 'react';
import { TextField, Button, Paper } from '@mui/material';
import { useAuth } from '../hooks/useAuth';

const LoginForm: React.FC = () => {
  const [credentials, setCredentials] = useState({
    email: '',
    password: ''
  });
  const { login } = useAuth();

  const handleSubmit = async (e: React.FormEvent) => {
    e.preventDefault();
    await login(credentials);
  };

  return (
    <Paper elevation={3} sx={{ p: 4 }}>
      <form onSubmit={handleSubmit}>
        <TextField
          fullWidth
          label="Email"
          value={credentials.email}
          onChange={(e) => setCredentials({
            ...credentials,
            email: e.target.value
          })}
        />
        <Button type="submit" variant="contained">
          Se connecter
        </Button>
      </form>
    </Paper>
  );
};
\end{lstlisting}

\section{Sécurité}

\subsection{Authentification}
\begin{itemize}
    \item \textbf{Hachage des mots de passe} : Utilisation de bcrypt avec salt
    \item \textbf{Tokens JWT} : Authentification stateless sécurisée
    \item \textbf{Middleware d'authentification} : Protection des routes sensibles
\end{itemize}

\subsection{Validation des Données}
\begin{itemize}
    \item Validation côté client avec des schémas TypeScript
    \item Validation côté serveur avec Mongoose
    \item Sanitisation des entrées utilisateur
\end{itemize}

\subsection{Protection CORS}
Configuration CORS restrictive pour limiter les accès depuis des domaines non autorisés.

\section{Tests}

\subsection{Tests Unitaires}
\begin{itemize}
    \item \textbf{Frontend} : Jest et React Testing Library
    \item \textbf{Backend} : Jest et Supertest pour les tests d'API
\end{itemize}

\subsection{Tests d'Intégration}
\begin{itemize}
    \item Tests des flux complets utilisateur
    \item Tests de l'API avec une base de données de test
\end{itemize}

\subsection{Tests End-to-End}
\begin{itemize}
    \item Cypress pour les tests d'interface utilisateur
    \item Scénarios de test complets : inscription, connexion, création de projet
\end{itemize}

\section{Déploiement}

\subsection{Environnements}
\begin{itemize}
    \item \textbf{Développement} : Serveurs locaux avec hot reload
    \item \textbf{Test} : Environnement de staging avec données de test
    \item \textbf{Production} : Déploiement sur services cloud
\end{itemize}

\subsection{CI/CD}
Pipeline automatisé incluant :
\begin{enumerate}
    \item Exécution des tests
    \item Build des applications
    \item Déploiement automatique
    \item Vérifications post-déploiement
\end{enumerate}

\section{Performance et Optimisation}

\subsection{Frontend}
\begin{itemize}
    \item Code splitting avec React.lazy
    \item Optimisation des images
    \item Mise en cache des requêtes API
    \item Progressive Web App (PWA)
\end{itemize}

\subsection{Backend}
\begin{itemize}
    \item Indexation des champs MongoDB fréquemment interrogés
    \item Pagination des résultats
    \item Mise en cache Redis pour les données fréquentes
    \item Compression gzip des réponses
\end{itemize}

\section{Monitoring et Maintenance}

\subsection{Logging}
\begin{itemize}
    \item Winston pour la journalisation côté serveur
    \item Niveaux de log appropriés (error, warn, info, debug)
    \item Rotation des fichiers de log
\end{itemize}

\subsection{Monitoring}
\begin{itemize}
    \item Surveillance des performances applicatives
    \item Monitoring de la base de données
    \item Alertes en cas de dysfonctionnement
\end{itemize}

\section{Évolutions Futures}

\subsection{Fonctionnalités Prévues}
\begin{itemize}
    \item Système de notation et d'avis
    \item Chat en temps réel entre freelancers et clients
    \item Système de paiement intégré
    \item Application mobile native
\end{itemize}

\subsection{Améliorations Techniques}
\begin{itemize}
    \item Migration vers Next.js pour le SSR
    \item Implémentation de GraphQL
    \item Microservices pour une meilleure scalabilité
    \item Intelligence artificielle pour le matching
\end{itemize}

\section{Conclusion}

Le projet LanceJob représente une solution complète et moderne pour la mise en relation entre freelancers et clients. L'architecture choisie assure une séparation claire des responsabilités, une sécurité robuste et une expérience utilisateur optimale.

Les technologies utilisées garantissent la scalabilité et la maintenabilité du système, tandis que les pratiques de développement adoptées (tests, CI/CD, monitoring) assurent la qualité et la fiabilité de la plateforme.

\subsection{Retour d'Expérience}
Ce projet a permis de mettre en pratique les meilleures pratiques du développement web moderne, de la conception à la mise en production, en passant par les tests et le déploiement.

\end{document}
